\section{Supported topology retains its own direction and domains}
\label{sec:supported}
For a finite context with column family $F$, let $K_F$ have the nonempty
faces $U$ contained in some actual column of $F$, and include the empty
simplex. Define the supported concept poset
\[
 \mathcal C_F=\{B\in\Meet(F): B\ne\varnothing,
                         \ \exists C\in F\ (B\subseteq C)\}.
\]
This is not obtained by deleting lattice extrema indiscriminately. The
nonempty face poset $\mathcal F(K_F)$ retracts onto $\mathcal C_F$ by
$U\mapsto\cl_F(U)$: the inclusion $s$ and closure $r$ satisfy
$rs=1$ and $1\leq sr$. Their order complexes are therefore homotopy
equivalent. This is the established supported FCA--Dowker bridge
\cite{Common,GP,BS}.

\begin{proposition}[Growth and partial operational transport]
For generated doctrines $G\subseteq G^+$, the forward comparisons are
$K_G\subseteq K_{G^+}$ and $\mathcal C_G\subseteq\mathcal C_{G^+}$.
On old faces the closures satisfy $\cl_{G^+}(U)\subseteq\cl_G(U)$;
this supplies an order homotopy, not an equality of the two closure routes.
For a shared program $p:i\to j$, its maximal supported domain is
\[
 K_p=\{\varnothing\}\cup
 \{\varnothing\ne U\subseteq X_i:X(p)[U]\in K_{G_j}\}.
\]
It is a subcomplex of $K_{G_i}$, and growth enlarges this domain.
For composable programs $p,q$ one has exactly
\[
 K_{qp}=\{U\in K_p:X(p)[U]\in K_q\}.
\]
The corresponding concept domain consists of nonempty $B\in M_i$ for
which $p_!(B)$ is supported. Closure intertwines these action spans,
including their domains. Backward projection $\rho^M$ is supported
only on new supported extents having an old common-test witness.
\end{proposition}
\begin{proof}
Old columns remain available and old extents remain closed, proving the
forward inclusions. New closure is smaller; both sides are nonempty and
have an old common witness, so they lie in the new supported poset.
An inequality of monotone maps induces the asserted order homotopy.
If $X(p)[U]\subseteq C\in G_j$, then
$U\subseteq X(p)^{-1}C\in G_i$, proving domain inclusion and also
showing that closing $U$ does not leave its domain. Equation (CN) gives
strict commutation with the concept action. If $X(qp)[U]$ has a
common column, pulling that column back along $q$ shows that $X(p)[U]$
already has a common column. This proves the domain composition formula;
the converse follows immediately. The concept version uses the same
argument and $q_!p_!=(qp)_!$. Finally $\cl_G(B)$ is old-supported
if and only if $B$ is contained in an old column. This is the exact
domain of the backward supported comparison.
\end{proof}
Growth yields commuting simplicial action spans and lax concept-action
inclusions. It need not yield partial maps with identical domains: an old
face may acquire a newly supported image. Nor does a partial action give
a global homology endomorphism. The rich-to-poor quotient and lattice
projection must not be confused with the poor-to-rich nerve inclusion.

\paragraph{A different mixed square can be strict.}
Evidence restriction is also distinct from language growth. If $E\subseteq F$
are columns on $H$ and $U\subseteq H$ is retained evidence, then for an
$F$-closed extent $A$,
\[
 \cl_E(A\cap U)\cap U\ \subseteq\ \cl_E(A)\cap U.
\]
Take $H=\{a,b,c\}$, $U=\{a,b\}$, old columns $\{a\},H$, and new
column $A=\{a,c\}$. The two sides are $\{a\}$ and $\{a,b\}$.
Both are supported and nonempty. If $U$ contains a representative of
every $F$-row, equality holds for every $F$-closed $A$: each member of
$A$ has a row-equivalent representative in $A\cap U$, so the two sets
have the same old common tests. This mixed right/left-adjoint comparison
is separate from the strict backward naturality proved above.
